\chapter{استراتژی‌های جست‌وجوی فایل}

هنگام کاوش در یک توزیع لینوکسی، بی‌درنگ با واقعیتی انکارناپذیر روبه‌رو می‌شوید: تعداد فایل‌های موجود در سیستم به‌مراتب فراتر از تصور اولیه است. یک نصب استاندارد لینوکس به‌راحتی شامل هزاران دایرکتوری و فایل است و همین امر، پرسشی بنیادین را پیش روی کاربر می‌گذارد: «چگونه می‌توان فایلی خاص را در این دریای اطلاعات یافت؟»

هرچند ساختار سیستم فایل لینوکس بر اساس قراردادهای ریشه‌دار در سیستم‌های شبه‌یونیکس به‌دقت سازمان‌دهی شده است، اما وسعت این ساختار سلسله‌مراتبی همچنان می‌تواند برای کاربر چالش‌برانگیز باشد. در این فصل، به بررسی دو ابزار بنیادین برای مکان‌یابی فایل‌ها خواهیم پرداخت:

\begin{itemize}
  \item \cmd{locate}: ابزاری سریع برای یافتن فایل‌ها بر اساس نام آن‌ها که از یک پایگاه‌دادهٔ از پیش ساخته‌شده استفاده می‌کند.
  \item \cmd{find}: ابزاری قدرتمند و انعطاف‌پذیر برای جست‌وجوی دقیق فایل‌ها که به‌صورت بلادرنگ (\en{Real-time}) سلسله‌مراتب دایرکتوری‌ها را پیمایش می‌کند.
\end{itemize}

علاوه بر این، دستوری را معرفی می‌کنیم که مکملی حیاتی برای ابزارهای جست‌وجو محسوب می‌شود و جهت پردازش فهرست فایل‌های یافت شده به‌کار می‌رود:

\begin{itemize}
  \item \cmd{xargs}: ابزاری برای ساخت و اجرای فرامین از طریق ورودی استاندارد (\en{Standard Input}). این دستور آیتم‌های ورودی را خوانده و آن‌ها را به‌عنوان آرگومان به فرمان مورد نظر منتقل می‌کند.
\end{itemize}

در نهایت، برای تکمیل مهارت‌های مدیریتی شما، دو دستور کاربردی دیگر را نیز بررسی خواهیم کرد:

\begin{itemize}
  \item \cmd{touch}: جهت تغییر برچسب‌های زمانی (\en{Timestamps}) مانند زمان دسترسی و زمان تغییر فایل. در صورت عدم وجود فایل، یک فایل خالی ایجاد می‌کند.
  \item \cmd{stat}: برای نمایش جزئیات دقیق و وضعیت آماری یک فایل یا سیستم فایل، شامل اطلاعاتی نظیر اندازه، مجوزها، تاریخ ایجاد و تاریخ دسترسی.
\end{itemize}

\section{جست‌وجوی سریع فایل‌ها با دستور locate}

ابزار \cmd{locate} فرآیند جست‌وجو را از طریق بررسی در یک پایگاه‌داده پیش‌ساخته از مسیر فایل‌ها (\en{Pathnames}) انجام می‌دهد که منجر به سرعتی بسیار بالا در بازیابی نتایج می‌شود. این برنامه هر مسیری را که شامل زیررشته (\en{Substring}) مورد نظر شما باشد، استخراج کرده و نمایش می‌دهد.

به‌عنوان مثال، اگر بخواهیم تمامی برنامه‌هایی که نام آن‌ها با عبارت \cmd{zip} آغاز می‌شود را بیابیم، با فرض اینکه فایل‌های اجرایی معمولاً در دایرکتوری‌هایی با پسوند \file{/bin} قرار دارند، می‌توانیم دستور را به شکل زیر اجرا کنیم:

\begin{latin}
\begin{lstlisting}
[me@linuxbox ~]$ locate bin/zip
\end{lstlisting}
\end{latin}

در این حالت، \cmd{locate} پایگاه‌داده را پیمایش کرده و هر مسیری را که شامل توالی کاراکتری \file{bin/zip} باشد، در خروجی لیست می‌کند:

\begin{latin}
\begin{lstlisting}
/usr/bin/zip
/usr/bin/zipcloak
/usr/bin/zipgrep
/usr/bin/zipinfo
/usr/bin/zipnote
/usr/bin/zipsplit
\end{lstlisting}
\end{latin}

برای جست‌وجوهای پیچیده‌تر، می‌توان توانمندی \cmd{locate} را با ابزارهایی نظیر \cmd{grep} ترکیب کرد تا نتایج دقیق‌تری به دست آید:

\begin{latin}
\begin{lstlisting}
[me@linuxbox ~]$ locate zip | grep bin
/bin/bunzip2
/bin/bzip2
/bin/bzip2recover
/bin/gunzip
/bin/gzip
/usr/bin/funzip
/usr/bin/gpg-zip
/usr/bin/preunzip
/usr/bin/prezip
/usr/bin/prezip-bin
/usr/bin/unzip
/usr/bin/unzipsfx
/usr/bin/zip
/usr/bin/zipcloak
/usr/bin/zipgrep
/usr/bin/zipinfo
/usr/bin/zipnote
/usr/bin/zipsplit
\end{lstlisting}
\end{latin}

ابزار \cmd{locate} پیشینه‌ای طولانی در اکوسیستم شبه‌یونیکس دارد و امروزه نسخه‌های بهینه‌شده‌ای از آن مانند \cmd{plocate} و \cmd{mlocate} در توزیع‌های مدرن لینوکس به کار گرفته می‌شوند؛ هرچند اغلب آن‌ها از طریق یک پیوند نمادین (\en{Symbolic Link}) و با همان نام استاندارد \cmd{locate} قابل فراخوانی هستند.

این نسخه‌ها اگرچه در گزینه‌های اصلی مشترک هستند، اما برخی از آن‌ها قابلیت‌های پیشرفته‌ای نظیر پشتیبانی از عبارت‌های باقاعده (\en{Regular Expressions}) و تطبیق با نویسه‌های جانشین (\en{Wildcards}) را نیز ارائه می‌دهند. جهت تشخیص نسخه دقیق نصب‌شده روی سیستم و مطالعه قابلیت‌های اختصاصی آن، می‌توانید به صفحه‌ی راهنمای برنامه (\cmd{man locate}) مراجعه کنید.

\section{سازوکار به‌روزرسانی پایگاه‌داده locate}

شاید مشاهده کرده باشید که در برخی توزیع‌ها، دستور \cmd{locate} بلافاصله پس از نصب سیستم‌عامل عمل نمی‌کند، اما پس از گذشت یک روز، به‌درستی نتایج را نمایش می‌دهد. دلیل این پدیده به ماهیت پایگاه‌داده‌ی این ابزار بازمی‌گردد.

پایگاه‌داده‌ای که \cmd{locate} از آن پرس‌وجو می‌کند، توسط برنامه‌ای به نام \cmd{updatedb} ایجاد و نگهداری می‌شود. این برنامه معمولاً به‌صورت دوره‌ای و از طریق یک وظیفه‌ی زمان‌بندی‌شده (\en{cron job}) اجرا می‌گردد. در اغلب توزیع‌های لینوکس، فرآیند اجرای \cmd{updatedb} به‌صورت پیش‌فرض یک‌بار در شبانه‌روز تنظیم شده است.

از آنجا که این پایگاه‌داده به‌صورت بلادرنگ (\en{Real-time}) به‌روزرسانی نمی‌شود، فایل‌هایی که در فواصل زمانی بین دو اجرای \cmd{updatedb} ایجاد شده‌اند، در نتایج جست‌وجوی \cmd{locate} ظاهر نخواهند شد.

برای رفع این محدودیت و اطمینان از وجود آخرین تغییرات در پایگاه‌داده، می‌توانید دستور \cmd{updatedb} را به‌صورت دستی و با امتیازات ابرکاربر (\en{superuser}) اجرا کنید:

\begin{latin}
\begin{lstlisting}
[me@linuxbox ~]$ sudo updatedb
\end{lstlisting}
\end{latin}

\section{پیمایش تخصصی سیستم فایل با find}

در حالی که ابزار \cmd{locate} تنها بر اساس نام فایل و با تکیه بر پایگاه‌داده عمل می‌کند، برنامه \cmd{find} با رویکردی مستقیم، یک دایرکتوری مشخص و تمامی زیرشاخه‌های آن را برای یافتن فایل‌ها بر اساس مجموعه‌ای از صفات و ویژگی‌های فنی پیمایش می‌کند. ما زمان قابل‌توجهی را صرف یادگیری \cmd{find} خواهیم کرد؛ چرا که قابلیت‌های پیشرفته‌ی آن در مباحث اسکریپت‌نویسی و مدیریت سیستم، کاربردهای فراوانی دارد.

در ساده‌ترین سناریوی اجرایی، \cmd{find} به همراه یک یا چند مسیر (\en{Path}) برای جست‌وجو فراخوانی می‌شود. برای مثال، جهت استخراج فهرست کامل محتویات دایرکتوری خانگی، از دستور زیر استفاده می‌کنیم:

\begin{latin}
\begin{lstlisting}
[me@linuxbox ~]$ find ~
\end{lstlisting}
\end{latin}

در حساب‌های کاربری فعال، این دستور حجم عظیمی از داده را تولید می‌کند. از آنجا که خروجی این ابزار به خروجی استاندارد (\en{stdout}) هدایت می‌شود، می‌توان آن را برای پردازش‌های بعدی به فرامین دیگر پایپ (\en{Pipe}) کرد. برای نمونه، جهت شمارش تعداد کل فایل‌ها و دایرکتوری‌ها، از ابزار \cmd{wc} بهره می‌گیریم:

\begin{latin}
\begin{lstlisting}
[me@linuxbox ~]$ find ~ | wc -l
47068
\end{lstlisting}
\end{latin}

قدرت واقعی \cmd{find} در توانایی تفکیک و شناسایی فایل‌ها بر اساس معیارهای دقیق نهفته است. این ابزار از ساختار ویژه‌ای شامل گزینه‌ها (\en{Options})، شرط‌ها (\en{Tests}) و کنش‌ها (\en{Actions}) استفاده می‌کند. برای درک عمیق این ابزار، ابتدا به بررسی مفهوم «شرط‌ها» می‌پردازیم که وظیفه‌ی فیلتر کردن نتایج را بر عهده دارند.

\subsection{فیلترهای جست‌وجو (Tests) در find}

تصور کنید قصد داریم خروجی جست‌وجو را منحصراً به دایرکتوری‌ها محدود کنیم؛ برای این منظور می‌توانیم از پارامتر تعیین نوع یا همان «شرط» \cmd{-type} استفاده کنیم:

\begin{latin}
\begin{lstlisting}
[me@linuxbox ~]$ find ~ -type d | wc -l
1695
\end{lstlisting}
\end{latin}

افزودن عبارت \cmd{-type d} فرآیند جست‌وجو را به‌گونه‌ای فیلتر می‌کند که تنها دایرکتوری‌ها (\en{Directories}) در فهرست نهایی لحاظ شوند. در مقابل، اگر بخواهیم جست‌وجو را صرفاً به فایل‌های معمولی (\en{Regular Files}) محدود نماییم، از آرگومان \cmd{-type f} استفاده می‌کنیم:

\begin{latin}
\begin{lstlisting}
[me@linuxbox ~]$ find ~ -type f | wc -l
38737
\end{lstlisting}
\end{latin}

\begin{table}[htbp]
\centering
\caption{انواع فایل در \cmd{find}}
\label{tab:find-file-types}
\begin{tabularx}{\textwidth}{c X}
\toprule
\textbf{شناسه} & \textbf{شرح فنی} \\
\midrule
\cmd{b} & فایل دستگاهی بلوکی (\en{Block Special Device}) \\
\addlinespace
\cmd{c} & فایل دستگاهی کاراکتری (\en{Character Special Device}) \\
\addlinespace
\cmd{d} & دایرکتوری (\en{Directory}) \\
\addlinespace
\cmd{f} & فایل معمولی (\en{Regular File}) \\
\addlinespace
\cmd{l} & پیوند نمادین (\en{Symbolic Link}) \\
\bottomrule
\end{tabularx}
\end{table}

علاوه بر نوع فایل، امکان جست‌وجو بر اساس مؤلفه‌هایی نظیر نام و حجم نیز فراهم است. برای مثال، جهت یافتن و شمارش تمامی فایل‌های معمولی که با الگوی \cmd{*.JPG} مطابقت دارند و حجم آن‌ها فراتر از یک مگابایت است، دستور زیر را به کار می‌بریم:

\begin{latin}
\begin{lstlisting}
[me@linuxbox ~]$ find ~ -type f -name "*.JPG" -size +1M | wc -l
840
\end{lstlisting}
\end{latin}

در این نمونه، شرط \cmd{-name} را به همراه یک نویسه جانشین (\en{Wildcard}) به‌کار برده‌ایم. شایان ذکر است که الگو حتماً باید درون کوتیشن قرار گیرد تا از تفسیر و بسط مسیر (\en{Pathname Expansion}) توسط پوسته پیش از اجرای دستور جلوگیری شود. در پروتکل \cmd{find}، هنگام استفاده از شرط \cmd{-size} به همراه یک مقدار عددی، علامت‌های \cmd{+} و \cmd{-} به‌ترتیب بیانگر «بزرگ‌تر از» و «کوچک‌تر از» آن مقدار هستند؛ در صورتی که هیچ‌یک از این دو علامت به کار نرود، جستجو دقیقاً بر روی همان مقدار اعمال خواهد شد.

\begin{table}[htbp]
\centering
\caption{واحدهای سنجش حجم در ابزار \cmd{find}}
\label{tab:find-size-units}
\begin{tabularx}{\textwidth}{c l X}
\toprule
\textbf{نویسه} & \textbf{واحد اندازه‌گیری} & \textbf{توضیحات فنی} \\
\midrule
\cmd{b} & بلوک‌های ۵۱۲ بایتی (\en{512-byte Blocks}) & واحد پیش‌فرض \cmd{find} در صورت عدم تعیین مقیاس؛ هر بلوک معادل ۵۱۲ بایت است. \\
\addlinespace
\cmd{c} & بایت (\en{Bytes}) & برای جست‌وجوی دقیق در سطح بایت؛ مناسب برای فایل‌های بسیار کوچک یا بررسی‌های دقیق. \\
\addlinespace
\cmd{w} & واژه‌های ۲ بایتی (\en{2-byte Words}) & هر واژه معادل ۲ بایت است؛ کاربرد آن عمدتاً در معماری‌های خاص و سیستم‌های قدیمی‌تر دیده می‌شود. \\
\addlinespace
\cmd{k} & کیبی‌بایت (\en{Kibibytes}) & معادل ۱,۰۲۴ بایت؛ رایج‌ترین واحد برای فایل‌های با اندازه‌ی متوسط. \\
\addlinespace
\cmd{M} & مبی‌بایت (\en{Mebibytes}) & معادل ۱,۰۴۸,۵۷۶ بایت (یعنی ۱۰۲۴ کیلوبایت)؛ مناسب برای فایل‌های حجیم مانند تصاویر و ویدئوها. \\
\addlinespace
\cmd{G} & گیبی‌بایت (\en{Gibibytes}) & معادل ۱,۰۷۳,۷۴۱,۸۲۴ بایت (یعنی ۱۰۲۴ مگابایت)؛ برای فایل‌های بسیار بزرگ مانند ایمیج‌های دیسک و پایگاه‌های داده. \\
\bottomrule
\end{tabularx}
\end{table}

ابزار \cmd{find} از طیف گسترده‌ای از شرط‌ها برای فیلتر کردن دقیق نتایج پشتیبانی می‌کند. جدول زیر نمایی کلی از پرکاربردترینِ این شرط‌ها را ارائه می‌دهد. شایان ذکر است در تمامی شرط‌هایی که مستلزم دریافت یک آرگومان عددی هستند، می‌توان از پیشوندهای \cmd{+} (بیشتر از) و \cmd{-} (کمتر از) برای تعیین بازه‌های مقداری استفاده کرد.

\begin{longtable}{l p{0.72\textwidth}}
\caption{پارامترهای شرطی (\en{Tests}) در دستور \cmd{find}}
\label{tab:find-tests} \\
\toprule
\textbf{شرط (\en{Test})} & \textbf{شرح عملکرد فنی} \\
\midrule
\endfirsthead

\multicolumn{2}{c}{\tablename\ \thetable{} (ادامه) \rule[-0.4em]{0pt}{1.4em}} \\
\toprule
\textbf{شرط (\en{Test})} & \textbf{شرح عملکرد فنی} \\
\midrule
\endhead

\midrule
\multicolumn{2}{r}{\footnotesize ادامه در صفحه‌ی بعد \ldots} \\
\endfoot

\bottomrule
\endlastfoot

\cmd{-cmin n} & تطبیق با فایل‌هایی که فرادادهٔ آن‌ها (مانند مجوزها یا مالکیت) دقیقاً \en{n} دقیقه پیش تغییر کرده است (\en{Status change}). برای تعیین بازه‌های زمانی دقیق‌تر، می‌توان از پیشوندهای \cmd{+} و \cmd{-} استفاده کرد: \cmd{-n} به معنای «کمتر از \en{n} دقیقه پیش»، \cmd{+n} به معنای «بیش از \en{n} دقیقه پیش»، و عدم استفاده از پیشوند به معنای «دقیقاً \en{n} دقیقه پیش» است. \\
\addlinespace
\cmd{-cnewer file} & یافتن فایل‌هایی که فرادادهٔ آن‌ها پس از آخرین تغییر فایل مرجع (\file{file}) تغییر کرده است. \\
\addlinespace
\cmd{-ctime n} & تطبیق با فایل‌هایی که فرادادهٔ آن‌ها (مانند مجوزها یا مالکیت) دقیقاً \en{n} روز پیش (هر روز به عنوان یک دورهٔ کامل ۲۴ ساعته محاسبه می‌شود) تغییر کرده است. مقدار \cmd{0} به معنای «کمتر از ۲۴ ساعت پیش»، \cmd{1} به معنای «بین ۲۴ تا ۴۸ ساعت پیش» و \cmd{+1} به معنای «بیش از ۴۸ ساعت پیش» است. \\
\addlinespace
\cmd{-empty} & شناسایی فایل‌ها و دایرکتوری‌های تهی (فایل‌های با حجم صفر و دایرکتوری‌های بدون محتوا). \\
\addlinespace
\cmd{-group name} & فیلتر کردن بر اساس گروه کاربری (نام گروه یا \en{GID} عددی). \\
\addlinespace
\cmd{-iname pattern} & مشابه \cmd{-name} اما بدون حساسیت به حروف کوچک و بزرگ (\en{Case-insensitive}). \\
\addlinespace
\cmd{-inum n} & جست‌وجو بر اساس شمارهٔ آی‌نود (\en{inode})؛ ابزاری کلیدی برای یافتن تمام پیوندهای سخت (\en{Hard Links}) یک فایل. \\
\addlinespace
\cmd{-mmin n} & تطبیق با فایل‌هایی که محتوای آن‌ها دقیقاً \en{n} دقیقه پیش اصلاح (\en{Modify}) شده است. \\
\addlinespace
\cmd{-mtime n} & تطبیق با فایل‌هایی که محتوای آن‌ها دقیقاً \en{n} روز پیش اصلاح شده است (\en{Modify time}). دقت کنید که منظور از «روز» در اینجا یک دورهٔ کامل ۲۴ ساعته است. برای تعیین بازه‌های زمانی دقیق‌تر، می‌توان از پیشوندهای \cmd{+} و \cmd{-} استفاده کرد: \cmd{-n} به معنای «کمتر از \en{n} روز پیش»، \cmd{+n} به معنای «بیش از \en{n} روز پیش»، و عدم استفاده از پیشوند به معنای «دقیقاً \en{n} روز پیش» است. \\
\addlinespace
\cmd{-name pattern} & تطبیق با نام فایل بر اساس نویسه‌های جانشین (\en{Wildcard}) مانند \cmd{*} و \cmd{?}. \\
\addlinespace
\cmd{-newer file} & یافتن فایل‌هایی که پس از فایل مرجع (\file{file}) اصلاح شده‌اند. \\
\addlinespace
\cmd{-nouser} & یافتن فایل‌های متعلق به کاربرانی که دیگر در سیستم وجود ندارند (مثلاً پس از حذف یک حساب کاربری). \\
\addlinespace
\cmd{-nogroup} & یافتن فایل‌های متعلق به گروه‌های نامعتبر یا حذف‌شده (گروهی که در \file{/etc/group} وجود ندارد). \\
\addlinespace
\cmd{-perm mode} & تطبیق با مجوزهای دسترسی مشخص (به صورت هشت‌تایی (\en{Octal}) مانند \cmd{755} یا نمادین (\en{Symbolic}) مانند \cmd{u=rwx,go=rx}). \\
\addlinespace
\cmd{-samefile name} & مشابه \cmd{-inum}؛ یافتن فایل‌هایی که شماره آی‌نود (\en{inode}) یکسانی با فایل هدف دارند. \\
\addlinespace
\cmd{-size n} & فیلتر کردن فایل‌ها بر اساس حجم دقیق (یا بازه‌ای با استفاده از \cmd{+} و \cmd{-}). \\
\addlinespace
\cmd{-type c} & جست‌وجو بر اساس نوع فایل (مانند \cmd{f} برای فایل معمولی و \cmd{d} برای دایرکتوری). \\
\addlinespace
\cmd{-user name} & فیلتر کردن بر اساس نام کاربری یا شناسهٔ عددی مالک (\en{UID}). \\

\end{longtable}

فهرست فوق تنها شامل موارد پرکاربرد است. برای دسترسی به جزئیات فنی دقیق‌تر و پارامترهای پیشرفته‌تر، می‌توانید به مستندات سیستمی این برنامه (\cmd{man find}) مراجعه کنید.

\subsection{عملگرهای منطقی (Operators) در find}

با وجود تنوع گسترده‌ی شرط‌ها در ابزار \cmd{find}، گاهی برای توصیف دقیق شرایط جست‌وجو، به برقراری روابط منطقی میان آن‌ها نیاز داریم. به عنوان مثال، تصور کنید قصد دارید امنیت فایل‌ها و زیرشاخه‌های دایرکتوری خانگی خود را بررسی کنید؛ در این سناریو باید به دنبال تمامی فایل‌هایی باشید که سطح دسترسی آن‌ها \cmd{0600} نیست و تمامی دایرکتوری‌هایی که سطح دسترسی‌شان با \cmd{0700} مطابقت ندارد.

خوشبختانه \cmd{find} امکان ترکیب شرط‌ها را از طریق عملگرهای منطقی فراهم آورده است تا ساخت عبارات شرطی پیچیده میسر شود. برای اجرای جست‌وجوی فوق، دستور را به شکل زیر تنظیم می‌کنیم:

\begin{latin}
\begin{lstlisting}
[me@linuxbox ~]$ find ~ \( -type f -not -perm 0600 \) -or \( -type d -not -perm 0700 \)
\end{lstlisting}
\end{latin}

ظاهر این دستور ممکن است در ابتدا دشوار به نظر برسد، اما با درک منطق عملگرها، ساختار آن بسیار شفاف خواهد بود. جدول زیر عملگرهای منطقی مورد استفاده در \cmd{find} را معرفی می‌کند:

\begin{table}[htbp]
\centering
\caption{عملگرهای منطقی در \cmd{find}}
\label{tab:find-operators}
\begin{tabularx}{\textwidth}{l X}
\toprule
\textbf{عملگر} & \textbf{شرح عملکرد فنی} \\
\midrule
\cmd{-and} & تنها در صورتی نتیجه مثبت است که شرط‌های هر دو طرف عملگر برقرار باشند. این عملگر را می‌توان به صورت کوتاه‌شده \cmd{-a} نیز به کار برد. لازم به ذکر است که در صورت عدم درج هرگونه عملگر بین دو شرط، \cmd{find} به‌طور پیش‌فرض عملگر \en{AND} را اعمال می‌کند. \\
\addlinespace
\cmd{-or} & در صورتی که حداقل یکی از شرط‌های طرفین برقرار باشد، تطبیق صورت می‌گیرد. فرم کوتاه‌شده آن \cmd{-o} است. \\
\addlinespace
\cmd{-not} & این عملگر نتیجه‌ی شرطِ بلافاصله بعد از خود را معکوس می‌کند (نقیض منطقی). معادل آن علامت تعجب \cmd{!} است. \\
\addlinespace
\cmd{( )} & برای گروه‌بندی عبارات و ساخت منطق‌های چندلایه استفاده می‌شود. پرانتزها ترتیب ارزیابی (\en{Precedence}) را کنترل می‌کنند. به‌طور پیش‌فرض، \cmd{find} عبارات را از چپ به راست پردازش می‌کند، اما با پرانتزگذاری می‌توان این اولویت را تغییر داد یا خوانایی دستور را بهبود بخشید. \\
\bottomrule
\end{tabularx}
\end{table}

\begin{note}
\textbf{نکته:} از آنجا که پرانتزها در محیط پوسته (\en{Shell}) دارای معنای خاصی هستند، هنگام استفاده در دستور \cmd{find} باید حتماً با نویسه‌ی گریز (\cmd{\textbackslash}) مهار شوند (\en{Escape}) تا مستقیماً به برنامه‌ی \cmd{find} پاس داده شوند. همچنین رعایت فاصله (\en{Space}) در دو طرف پرانتز الزامی است؛ مانند:
\begin{latin}
\begin{lstlisting}
find ~ \( -type f \)
\end{lstlisting}
\end{latin}
\end{note}

با استناد به فهرست عملگرهای بررسی‌شده، اجازه دهید ساختار دستور \cmd{find} مورد نظرمان را کالبدشکافی کنیم. در یک نگاه کلی، مشاهده می‌کنیم که شرایط جست‌وجو در قالب دو گروه منطقی مجزا سازمان‌دهی شده‌اند که توسط عملگر \cmd{-or} به یکدیگر پیوند خورده‌اند:

\begin{latin}
\begin{lstlisting}
( expression 1 ) -or ( expression 2 )
\end{lstlisting}
\end{latin}

این ساختار کاملاً منطقی است؛ زیرا ما به‌طور هم‌زمان به‌دنبال فایل‌هایی با یک الگوی دسترسی خاص یا دایرکتوری‌هایی با الگوی دسترسی متفاوت هستیم. شاید این پرسش مطرح شود که چرا به‌جای \cmd{-and} از \cmd{-or} استفاده کرده‌ایم؟ پاسخ در ماهیت اشیاء سیستم فایل نهفته است: یک موجودیت در لحظه یا یک فایل معمولی است یا یک دایرکتوری، و نمی‌تواند هر دو باشد. بنابراین، برای پوشش هر دو حالت، باید از ترکیب زیر استفاده کنیم:

\begin{latin}
\begin{lstlisting}
( file with bad perms ) -or ( directory with bad perms )
\end{lstlisting}
\end{latin}

گام بعدی، تعریف دقیق «دسترسی غیرمجاز» است. به‌جای تلاش برای فیلتر کردن بی‌شمار حالت ناامن، ما از رویکرد نقیض استفاده می‌کنیم؛ یعنی مواردی را جست‌وجو می‌کنیم که با «دسترسی مجاز» مطابقت ندارند. با فرض اینکه سطح دسترسی بهینه برای فایل‌ها \cmd{0600} و برای دایرکتوری‌ها \cmd{0700} است، عبارت شرطی برای یافتن فایل‌های ناامن به شرح زیر خواهد بود:

\begin{latin}
\begin{lstlisting}
-type f -and -not -perm 0600
\end{lstlisting}
\end{latin}

و برای دایرکتوری‌های ناامن:

\begin{latin}
\begin{lstlisting}
-type d -and -not -perm 0700
\end{lstlisting}
\end{latin}

مطابق با آنچه در جدول فوق اشاره شد، عملگر \cmd{-and} به‌صورت ضمنی (\en{Implicit}) در نظر گرفته می‌شود و نیازی به درج صریح آن نیست. با تجمیع این مولفه‌ها، به ساختار نهایی زیر می‌رسیم:

\begin{latin}
\begin{lstlisting}
find ~ ( -type f -not -perm 0600 ) -or ( -type d -not -perm 0700 )
\end{lstlisting}
\end{latin}

در نهایت، به دلیل جایگاه ویژه پرانتزها در نحو (\en{Syntax}) پوسته، مهار کردن آن‌ها با نویسه گریز \cmd{\textbackslash} الزامی است تا مستقیماً به ابزار \cmd{find} ارسال شوند:

\begin{latin}
\begin{lstlisting}
find ~ \( -type f -not -perm 0600 \) -or \( -type d -not -perm 0700 \)
\end{lstlisting}
\end{latin}

ویژگی کلیدی دیگری در عملکرد عملگرهای منطقی وجود دارد که درک آن برای تسلط بر این ابزار ضروری است. سناریویی را در نظر بگیرید که در آن دو عبارت شرطی توسط یک عملگر به یکدیگر پیوند خورده‌اند:

\begin{latin}
\begin{lstlisting}
expr1 -operator expr2
\end{lstlisting}
\end{latin}

در تمامی حالات، ابتدا \en{expr1} مورد ارزیابی قرار می‌گیرد؛ اما تداوم فرآیند و اجرای \en{expr2} کاملاً به نوع عملگر و خروجی عبارت اول بستگی دارد. این سازوکار که در علوم کامپیوتر با عنوان ارزیابی اتصال‌کوتاه (\en{Short-circuit Evaluation}) شناخته می‌شود، در جدول زیر تشریح شده است:

\begin{table}[htbp]
\centering
\caption{منطق ارزیابی شرطی در \cmd{find}}
\label{tab:find-short-circuit}
\begin{tabularx}{\textwidth}{l l X}
\toprule
\textbf{نتیجه ارزیابی \en{Expr1}} & \textbf{عملگر منطقی} & \textbf{آیا \en{expr2} اجرا می‌شود؟} \\
\midrule
\en{True} (صحیح) & \cmd{-and} & بله \\
\addlinespace
\en{False} (غلط) & \cmd{-and} & خیر \\
\addlinespace
\en{True} (صحیح) & \cmd{-or} & خیر \\
\addlinespace
\en{False} (غلط) & \cmd{-or} & بله \\
\bottomrule
\end{tabularx}
\end{table}

دلیل اصلی پیاده‌سازی این منطق، بهینه‌سازی عملکرد (\en{Performance}) سیستم است. در عملگر \cmd{-and}، چنانچه \en{expr1} نادرست باشد، کلیت عبارت تحت هر شرایطی نادرست خواهد بود؛ لذا صرف زمان برای پردازش \en{expr2} منطقی نیست. به همین ترتیب در عملگر \cmd{-or}، اگر \en{expr1} صحیح ارزیابی شود، نتیجه نهاییِ عبارت بدون توجه به محتوای \en{expr2} صحیح خواهد بود.

علاوه بر جنبه‌ی سرعت، اهمیت حیاتی این موضوع زمانی آشکار می‌شود که قصد داریم اجرای کنش‌ها (\en{Actions}) را کنترل کنیم. از آنجا که در سیستم‌عامل لینوکس بسیاری از کنش‌ها تنها در صورت موفقیت‌آمیز بودن شرط‌های قبلی اجرا می‌شوند، درک دقیق این اولویت‌بندی به ما اجازه می‌دهد تا فرآیندهای پیچیده‌ی مدیریتی را با دقت بالایی خودکارسازی کنیم.

\subsection{کنش‌های از پیش تعریف‌شده (Predefined Actions)}

پس از یافتن فایل‌های مورد نظر، گام نهایی و عملیاتی، اتخاذ تصمیم و اجرای فرآیند بر روی نتایج به‌دست‌آمده است. ابزار \cmd{find} این امکان را فراهم می‌کند تا بر اساس خروجیِ فیلترها، عملیات‌های مشخصی را به‌صورت خودکار اجرا کنید. در لینوکس، مجموعه‌ای از کنش‌های از پیش تعریف‌شده (\en{Predefined Actions}) برای این منظور تعبیه شده است که کارهای رایجی مانند حذف یا نمایش جزئیات فایل‌ها را با یک پارامتر ساده ممکن می‌سازند.

با این حال، اگر کنش‌های از پیش تعریف‌شده پاسخگوی نیاز شما نباشند، \cmd{find} امکان فراخوانی دستورات سفارشی (\en{User-defined Actions}) را نیز فراهم می‌کند. با استفاده از پارامتر \cmd{-exec}، می‌توانید هر دستور دلخواهی را روی فایل‌های یافت‌شده اجرا کنید. این قابلیت، قدرت مانور بالایی به کاربر می‌دهد و امکان انجام عملیات‌های پیچیده و شخصی‌سازی‌شده را فراهم می‌آورد.

در جدول زیر، برخی از پرکاربردترینِ این کنش‌ها تشریح شده‌اند.

\begin{table}[htbp]
\centering
\caption{کنش‌های از پیش تعریف‌شده در ابزار \cmd{find}}
\label{tab:find-actions}
\begin{tabularx}{\textwidth}{l X}
\toprule
\textbf{کنش} & \textbf{شرح عملکرد فنی} \\
\midrule
\cmd{-delete} & حذف آنی فایل یا دایرکتوری منطبق با معیارهای جست‌وجو. این کنش باید با احتیاط کامل به‌کار رود؛ زیرا فایل‌ها را به‌صورت دائمی حذف می‌کند. \\
\addlinespace
\cmd{-ls} & نمایش جزئیات فایل منطبق، مشابه با خروجی دستور \cmd{ls -dils}. اطلاعات به خروجی استاندارد (\en{stdout}) ارسال می‌شود. \\
\addlinespace
\cmd{-print} & چاپ مسیر کامل فایل یافت‌شده. این کنش به‌صورت پیش‌فرض (\en{Default}) فعال است؛ یعنی اگر هیچ کنشی تعریف نکنید، \cmd{find} این دستور را اجرا می‌کند. \\
\addlinespace
\cmd{-quit} & توقف کامل فرآیند جست‌وجو بلافاصله پس از یافتن نخستین موردِ منطبق با معیارها. این کنش در سناریوهایی که تنها به یک نتیجه نیاز دارید (مانند یافتن سریع یک فایل خاص)، بسیار کارآمد است. \\
\bottomrule
\end{tabularx}
\end{table}

شایان ذکر است که این فهرست تنها بخشی از قابلیت‌های عملیاتی این ابزار است. برای بررسی تمامی کنش‌ها و جزئیات فنی دقیق‌تر، مطالعه‌ی مستندات سیستمی (\cmd{man find}) توصیه می‌شود.

نکته‌ی حیاتی در استفاده از کنش‌ها، به‌ویژه کنش‌های تخریبی مانند \cmd{-delete}، ترکیب آن‌ها با منطق ارزیابی اتصال‌کوتاه است. از آنجا که کنش‌ها نیز مانند شرط‌ها یک مقدار بازگشتی (\en{True/False}) دارند، ترتیب قرارگیری آن‌ها در خط فرمان تعیین می‌کند که آیا عملیات روی فایل اجرا شود یا خیر. در نخستین مثالی که بررسی کردیم، دستور زیر را اجرا نمودیم:

\begin{latin}
\begin{lstlisting}
find ~
\end{lstlisting}
\end{latin}

این دستور فهرستی جامع از تمامی فایل‌ها و زیرشاخه‌های موجود در دایرکتوری خانگی کاربر تولید کرد. دلیل مشاهده این خروجی آن است که در صورت عدم تعیین صریح یک عملیات، کنش \cmd{-print} به‌صورت ضمنی (\en{Implicit}) توسط \cmd{find} اعمال می‌شود. بنابراین، دستور فوق با عبارت زیر کاملاً هم‌ارز است:

\begin{latin}
\begin{lstlisting}
find ~ -print
\end{lstlisting}
\end{latin}

همچنین می‌توان از توانمندی \cmd{find} برای حذف فایل‌هایی که با معیارهای مشخصی مطابقت دارند، بهره جست. به‌طور مشخص، جهت پاکسازی فایل‌هایی با پسوند \file{.bak} (که غالباً به فایل‌های پشتیبان اطلاق می‌شود)، می‌توان دستور زیر را به کار برد:

\begin{latin}
\begin{lstlisting}
find ~ -type f -name '*.bak' -delete
\end{lstlisting}
\end{latin}

در این سناریو، تمامی محتویات دایرکتوری خانگی و زیرشاخه‌های آن برای یافتن فایل‌های معمولی که نامشان به پسوند \file{.bak} ختم می‌شود، پیمایش شده و بلافاصله پس از انطباق، حذف می‌گردند.

\begin{warning}
\textbf{هشدار امنیتی:} بدیهی است که هنگام استفاده از کنش \cmd{-delete} باید نهایت احتیاط را به کار بست. اکیداً توصیه می‌شود پیش از اجرای نهایی دستور، ابتدا آن را با جایگزین کردن کنش \cmd{-print} به‌جای \cmd{-delete} آزمایش کنید تا از صحت نتایج جست‌وجو و عدم حذف ناخواسته فایل‌ها اطمینان حاصل نمایید.
\end{warning}

پیش از ادامه، ضروری است تحلیل دقیق‌تری از تأثیر عملگرهای منطقی بر روند اجرای کنش‌ها داشته باشیم. دستور زیر را در نظر بگیرید:

\begin{latin}
\begin{lstlisting}
find ~ -type f -name '*.bak' -print
\end{lstlisting}
\end{latin}

همان‌طور که پیش‌تر اشاره شد، این دستور به دنبال فایل‌های معمولی (\cmd{-type f}) با پسوند \file{.bak} گشته و مسیر آن‌ها را در خروجی استاندارد چاپ می‌کند (\cmd{-print}). دلیل این رفتار، وجود روابط منطقیِ حاکم میان شرط‌ها و کنش‌ها است. به خاطر داشته باشید که به‌صورت پیش‌فرض، بین هر مؤلفه یک عملگر \cmd{-and} ضمنی برقرار است. برای درک بهتر، می‌توان دستور را بدین صورت بازنویسی کرد:

\begin{latin}
\begin{lstlisting}
find ~ -type f -and -name '*.bak' -and -print
\end{lstlisting}
\end{latin}

در این ساختار، منطق ارزیابی اتصال‌کوتاه (\en{Short-circuit Evaluation}) تعیین می‌کند که هر بخش چه زمانی فراخوانی شود:

\begin{table}[htbp]
\centering
\caption{ترتیب اجرای شرط‌ها و کنش‌ها در \cmd{find}}
\label{tab:find-exec-order}
\begin{tabularx}{\textwidth}{l X}
\toprule
\textbf{شرط یا کنش} & \textbf{شرط اجرا} \\
\midrule
\cmd{-print} & تنها در صورتی اجرا می‌شود که هر دو شرط قبلی صحیح باشند. \\
\addlinespace
\cmd{-name '*.bak'} & تنها در صورتی اجرا می‌شود که خروجی \cmd{-type f} صحیح باشد. \\
\addlinespace
\cmd{-type f} & همواره اجرا می‌شود (نخستین عبارت در زنجیره‌ی منطقی). \\
\bottomrule
\end{tabularx}
\end{table}

از آنجا که این روابط منطقی تعیین‌کننده‌ی توالی اجرا هستند، ترتیب قرارگیری شرط‌ها و کنش‌ها در دستور حیاتی است. برای مثال، اگر کنش \cmd{-print} را به ابتدای زنجیره منتقل کنیم، رفتار ابزار به‌کلی تغییر می‌کند:

\begin{latin}
\begin{lstlisting}
find ~ -print -and -type f -and -name '*.bak'
\end{lstlisting}
\end{latin}

در این نسخه، چون کنش \cmd{-print} همواره صحیح (\en{True}) ارزیابی می‌شود، مسیر تمامی فایل‌ها و دایرکتوری‌ها بدون استثنا چاپ شده و سپس شرط‌های نوع و نام فایل روی آن‌ها انجام می‌گیرد (که در این مرحله عملاً بی‌اثر هستند).

\subsection{فراخوانی دستورات سفارشی (User-Defined Actions)}

علاوه بر عملیات‌های پیش‌فرض، ابزار \cmd{find} امکان اجرای فرامین دلخواه را بر روی نتایج جست‌وجو فراهم می‌کند. متداول‌ترین شیوه برای این کار، استفاده از کنش \cmd{-exec} است که ساختار نحوی آن به شرح زیر است:

\begin{latin}
\begin{lstlisting}
-exec command {} ;
\end{lstlisting}
\end{latin}

در این الگو، \file{command} نام دستور مورد نظر، مجموعه‌ی \cmd{\{\}} نمادی برای جایگذاری مسیر فایلِ یافت‌شده، و علامت سمی‌کالن (\cmd{;}) جداکننده‌ای اجباری برای اتمام دستور است. برای نمونه، جهت شبیه‌سازی عملکرد کنش \cmd{-delete} با استفاده از دستور \cmd{rm}، ساختار زیر به کار می‌رود:

\begin{latin}
\begin{lstlisting}
-exec rm '{}' ';'
\end{lstlisting}
\end{latin}

به دلیل اینکه آکولادها و سمی‌کالن در محیط پوسته دارای معانی رزرو شده هستند، مهار کردن آن‌ها با استفاده از کوتیشن یا نویسه‌ی گریز (\cmd{\textbackslash}) جهت ارسال صحیح به برنامه‌ی \cmd{find} الزامی است.

علاوه بر این، برای اجرای تعاملی و تحت نظارتِ دستورات، می‌توان از کنش \cmd{-ok} به‌جای \cmd{-exec} استفاده کرد. در این حالت، سیستم پیش از اجرای دستور بر روی هر فایل، از کاربر تأییدیه می‌گیرد:

\begin{latin}
\begin{lstlisting}
find ~ -type f -name 'foo*' -ok ls -l '{}' ';'
< ls ... /home/me/bin/foo > ? y
-rw-r-xr-x 1 me   me 224 2007-10-29 18:44 /home/me/bin/foo
< ls ... /home/me/foo.txt > ? y
-rw-r--r-- 1 me   me   0 2016-09-19 12:53 /home/me/foo.txt
\end{lstlisting}
\end{latin}

در مثال فوق، تمامی فایل‌هایی که نامشان با رشته‌ی \en{foo} آغاز می‌شود شناسایی شده و دستور \cmd{ls -l} برای تک‌تک آن‌ها فراخوانی می‌گردد؛ با این تفاوت که به‌واسطه‌ی استفاده از \cmd{-ok}، اجرای هر مرحله منوط به پاسخ مثبت (\cmd{y}) کاربر است.

\section{بهینه‌سازی عملکرد فرآیند جست‌وجو}

در حالت استاندارد با استفاده از کنش \cmd{-exec}، به ازای هر فایلِ منطبق، یک نمونه (\en{instance}) جداگانه از دستور مورد نظر فراخوانی می‌گردد. در پروژه‌هایی با تعداد فایل‌های انبوه، این رویکرد به دلیل سربارِ (\en{overhead}) ایجاد فرآیندهای متعدد، کارایی سیستم را کاهش می‌دهد. برای درک بهتر، تفاوت در این است که به‌جای اجرای انفرادی دستورات:

\begin{latin}
\begin{lstlisting}
ls -l file1
ls -l file2
\end{lstlisting}
\end{latin}

تمامی نتایج تجمیع شده و دستور تنها یک‌بار به شکل زیر اجرا شود:

\begin{latin}
\begin{lstlisting}
ls -l file1 file2
\end{lstlisting}
\end{latin}

برای دستیابی به این سطح از بهینه‌سازی، دو متدولوژی وجود دارد: استفاده از ابزار خارجی \cmd{xargs} و یا بهره‌گیری از قابلیت داخلی و مدرن در خودِ \cmd{find}. ابتدا رویکرد داخلی را بررسی می‌کنیم.

با جایگزین کردن علامت سمی‌کالن (\cmd{;}) در انتهای دستور با علامت جمع (\cmd{+})، ابزار \cmd{find} تمامی نتایج جست‌وجو را در قالب یک لیستِ آرگومان واحد تجمیع کرده و دستور را تنها یک‌بار فراخوانی می‌کند.

در مثال پیشین، دستور زیر به ازای هر فایل پیدا شده، یک‌بار \cmd{ls} را اجرا می‌کرد:

\begin{latin}
\begin{lstlisting}
find ~ -type f -name 'foo*' -exec ls -l '{}' ';'
-rw-r-xr-x 1 me   me  224 2007-10-29 18:44 /home/me/bin/foo
-rw-r--r-- 1 me   me    0 2016-09-19 12:53 /home/me/foo.txt
\end{lstlisting}
\end{latin}

اما با تغییر نحو دستور به شکل زیر، اگرچه خروجی ظاهراً یکسان است، اما سیستم از نظر عملیاتی بسیار بهینه‌تر عمل کرده و تنها یک پروسه برای \cmd{ls} ایجاد می‌کند:

\begin{latin}
\begin{lstlisting}
find ~ -type f -name 'foo*' -exec ls -l '{}' +
-rw-r-xr-x 1 me   me  224 2007-10-29 18:44 /home/me/bin/foo
-rw-r--r-- 1 me   me    0 2016-09-19 12:53 /home/me/foo.txt
\end{lstlisting}
\end{latin}

\section{مدیریت آرگومان‌های ورودی با ابزار xargs}

فرمان \cmd{xargs} قابلیتی منحصربه‌فرد در زنجیره دستورات لینوکس ارائه می‌دهد؛ این ابزار داده‌ها را از ورودی استاندارد (\en{stdin}) دریافت کرده و آن‌ها را به لیست آرگومان‌های یک دستور مشخص تبدیل می‌کند. در سناریوی مورد بحث، نحوه به‌کارگیری آن به شرح زیر است:

\begin{latin}
\begin{lstlisting}
find ~ -type f -name 'foo*' -print | xargs ls -l
-rw-r-xr-x 1 me   me  224 2007-10-29 18:44 /home/me/bin/foo
-rw-r--r-- 1 me   me    0 2016-09-19 12:53 /home/me/foo.txt
\end{lstlisting}
\end{latin}

در این ساختار، خروجی متنیِ \cmd{find} از طریق پایپ (\en{pipe}) به \cmd{xargs} منتقل می‌شود. سپس \cmd{xargs} این مسیرها را به عنوان آرگومان‌های ورودی برای فرمان \cmd{ls} سازمان‌دهی کرده و آن را اجرا می‌نماید.

\begin{note}
\textbf{نکته:} اگرچه ظرفیت پذیرش آرگومان در خط فرمان بسیار وسیع است، اما این مقدار نامحدود نیست. در مواردی که حجم نتایج جست‌وجو بسیار بالا باشد، ممکن است طول رشته‌ی نهایی از حد مجاز سیستم عبور کرده و منجر به بروز خطا در پوسته شود. ابزار \cmd{xargs} با هوشمندی این چالش را مدیریت می‌کند؛ بدین صورت که اگر طول خط فرمان از حد آستانه فراتر رود، دستور مورد نظر را در چندین مرحله (با حداکثر آرگومان مجاز در هر نوبت) اجرا می‌کند تا تمامی ورودی‌ها پردازش شوند. جهت مشاهده سقف مجاز طول خط فرمان در سیستم خود، می‌توانید از دستور \cmd{xargs --show-limits} استفاده کنید.
\end{note}

\section{مدیریت نام‌های غیرمتعارف در سیستم فایل}

سیستم‌های شبه‌یونیکس (\en{Unix-like}) اجازه می‌دهند که نام فایل‌ها شامل فواصل خالی (\en{Space}) و حتی کاراکترهای خط جدید (\en{Newline}) باشد. این انعطاف‌پذیری برای ابزارهایی نظیر \cmd{xargs} که وظیفه‌ی ساخت لیست آرگومان‌ها را بر عهده دارند، چالش‌آفرین است؛ چرا که به‌طور پیش‌فرض، فواصل خالی به عنوان جداکننده‌ی آرگومان‌ها تفسیر می‌شوند. در چنین شرایطی، فایلی با نام \file{My File.jpg} توسط دستور نهایی نه به عنوان یک موجودیت واحد، بلکه به صورت دو آرگومان مجزای \file{My} و \file{File.jpg} در نظر گرفته می‌شود که قطعاً منجر به بروز خطا خواهد شد.

برای رفع این نقیصه، دستورات \cmd{find} و \cmd{xargs} امکان استفاده از کاراکتر \en{Null} را به عنوان جداکننده‌ی آرگومان‌ها فراهم کرده‌اند. کاراکتر \en{Null} در کدگذاری \en{ASCII} با مقدار عددی صفر تعریف شده است (برخلاف کاراکتر فاصله (\en{Space}) که مقدار آن در جدول کدهای \en{ASCII} برابر با ۳۲ است) و چون این کاراکتر نمی‌تواند بخشی از نام یک فایل باشد، جداکننده‌ای کاملاً ایمن محسوب می‌شود.

در ابزار \cmd{find}، کنش \cmd{-print0} خروجی را به جای خط جدید با کاراکتر \en{Null} جدا می‌کند. در مقابل، دستور \cmd{xargs} با گزینه \cmd{--null} (یا فرم کوتاه‌شده‌ی \cmd{-0}) ورودی‌های تفکیک‌شده با \en{Null} را می‌پذیرد. نمونه‌ی پیاده‌سازی این تکنیک به شرح زیر است:

\begin{latin}
\begin{lstlisting}
find ~ -iname '*.jpg' -print0 | xargs --null ls -l
\end{lstlisting}
\end{latin}

با بهره‌گیری از این متدولوژی، می‌توان اطمینان حاصل کرد که تمامی فایل‌ها، حتی مواردی که دارای نام‌های پیچیده یا شامل فواصل داخلی هستند، بدون هیچ‌گونه خطا در زنجیره‌ی پردازشی مدیریت می‌شوند.

\section{پیاده‌سازی سناریو و بازسازی محیط آزمایشگاهی}

اکنون زمان آن فرا رسیده است تا آموخته‌های خود پیرامون ابزار \cmd{find} را در یک سناریوی عملیاتی به کار بگیریم. ابتدا، یک محیط تمرینی وسیع شامل تعداد زیادی زیرشاخه و فایل ایجاد می‌کنیم:

\begin{latin}
\begin{lstlisting}
[me@linuxbox ~]$ mkdir -p playground/dir-{001..100}
[me@linuxbox ~]$ touch playground/dir-{001..100}/file-{A..Z}
\end{lstlisting}
\end{latin}

قدرت و بهره‌وری خط فرمان در اینجا کاملاً مشهود است؛ تنها با دو دستور، دایرکتوری \file{playground} را ایجاد کردیم که شامل ۱۰۰ زیرشاخه است و هر یک از آن‌ها ۲۶ فایل مجزا را در خود جای داده‌اند. اجرای چنین عملیاتی در رابط‌های گرافیکی (\en{GUI}) بسیار زمان‌بر و دشوار خواهد بود.

در این فرآیند، از ترکیب دستور آشنای \cmd{mkdir}، قابلیت بسط آکولاد \en{Brace Expansion} در پوسته و دستور کاربردی \cmd{touch} بهره بردیم. استفاده از گزینه \cmd{-p} در دستور \cmd{mkdir} تضمین می‌کند که تمامی دایرکتوری‌های میانی و بالادستی در مسیر تعیین‌شده ساخته شوند. دستور \cmd{touch} نیز به‌طور معمول برای به‌روزرسانی برچسب‌های زمانی (\en{Timestamps}) فایل‌ها به کار می‌رود، اما در صورتی که فایل هدف وجود نداشته باشد، یک فایل تهی با آن نام ایجاد می‌کند.

اکنون که محیط آماده است، برای یافتن تمامی ۱۰۰ نسخه‌ی فایل موسوم به \file{file-A} از دستور زیر استفاده می‌کنیم:

\begin{latin}
\begin{lstlisting}
[me@linuxbox ~]$ find playground -type f -name 'file-A'
\end{lstlisting}
\end{latin}

شایان ذکر است که برخلاف دستور \cmd{ls}، خروجی \cmd{find} لزوماً مرتب‌شده (\en{Sorted}) نیست؛ چرا که ترتیب نمایش نتایج، تابع نحوه‌ی ذخیره‌سازی و چیدمان داده‌ها در سیستم فایل (\en{Disk order}) است. برای اطمینان از صحت عملیات، می‌توانیم با استفاده از ابزار \cmd{wc} تعداد نتایج را شمارش کنیم:

\begin{latin}
\begin{lstlisting}
[me@linuxbox ~]$ find playground -type f -name 'file-A' | wc -l
100
\end{lstlisting}
\end{latin}

در این بخش به بررسی نحوه‌ی فیلتر کردن فایل‌ها بر اساس زمان آخرین تغییرات آن‌ها می‌پردازیم؛ مهارتی که در فرآیندهایی نظیر تهیه نسخه‌های پشتیبان (\en{Backup}) یا سازمان‌دهی زمانی داده‌ها بسیار حیاتی است. برای شروع، ابتدا یک فایل مرجع (\en{Reference File}) ایجاد می‌کنیم تا از زمان اصلاح آن به عنوان مبنای مقایسه استفاده کنیم:

\begin{latin}
\begin{lstlisting}
[me@linuxbox ~]$ touch playground/timestamp
\end{lstlisting}
\end{latin}

این دستور فایلی با نام \file{timestamp} ایجاد کرده و برچسب زمانی آن را با زمان فعلی سیستم همگام‌سازی می‌کند. برای مشاهده‌ی دقیق جزئیات این فایل، از دستور \cmd{stat} استفاده می‌کنیم. این ابزار در واقع نسخه‌ای پیشرفته‌تر از \cmd{ls} است که تمامی فراداده‌های فایل و ویژگی‌های ساختاری آن را نمایش می‌دهد:

\begin{latin}
\begin{lstlisting}
[me@linuxbox ~]$ stat playground/timestamp
File: `playground/timestamp'
Size:   0           Blocks: 0          IO Block: 4096 regular empty file
Device: 803h/2051d    Inode: 14265061    Links: 1
Access: (0644/-rw-r--r--)  Uid: ( 1001/ me)   Gid: ( 1001/ me)
Access: 2025-10-08 15:15:39.000000000 -0400
Modify: 2025-10-08 15:15:39.000000000 -0400
Change: 2025-10-08 15:15:39.000000000 -0400
\end{lstlisting}
\end{latin}

چنانچه مجدداً دستور \cmd{touch} را بر روی این فایل اجرا کنیم، با بررسی دوباره توسط \cmd{stat} مشاهده خواهیم کرد که مقادیر زمانی به‌روزرسانی شده‌اند:

\begin{latin}
\begin{lstlisting}
[me@linuxbox ~]$ touch playground/timestamp
[me@linuxbox ~]$ stat playground/timestamp
File: `playground/timestamp'
Size:   0           Blocks: 0          IO Block: 4096 regular empty file
Device: 803h/2051d    Inode: 14265061    Links: 1
Access: (0644/-rw-r--r--)  Uid: ( 1001/ me)   Gid: ( 1001/ me)
Access: 2025-10-08 15:23:33.000000000 -0400
Modify: 2025-10-08 15:23:33.000000000 -0400
Change: 2025-10-08 15:23:33.000000000 -0400
\end{lstlisting}
\end{latin}

اکنون از ابزار \cmd{find} برای به‌روزرسانی تعدادی از فایل‌های موجود در محیط آزمایشگاهی (\file{playground}) استفاده می‌کنیم:

\begin{latin}
\begin{lstlisting}
[me@linuxbox ~]$ find playground -type f -name 'file-B' -exec touch '{}' ';'
\end{lstlisting}
\end{latin}

این دستور تمامی فایل‌های موسوم به \file{file-B} را شناسایی کرده و برچسب زمانی آن‌ها را به‌روزرسانی می‌کند. حال با بهره‌گیری از قابلیت مقایسه‌ی زمانی، فایل‌هایی را که پس از ایجاد فایل مرجع \file{timestamp} تغییر یافته‌اند، شناسایی می‌کنیم:

\begin{latin}
\begin{lstlisting}
[me@linuxbox ~]$ find playground -type f -newer playground/timestamp
\end{lstlisting}
\end{latin}

خروجی شامل هر ۱۰۰ نسخه‌ی فایل \file{file-B} خواهد بود؛ زیرا از آنجا که عملیات \cmd{touch} روی این فایل‌ها پس از ایجاد فایل \file{timestamp} صورت گرفته است، آن‌ها اکنون نسبت به مرجع ما «جدیدتر» (\en{Newer}) محسوب می‌شوند.

در نهایت، بیایید شرط بررسی سطوح دسترسی را که پیش‌تر آموختیم، در محیط \file{playground} عملیاتی کنیم:

\begin{latin}
\begin{lstlisting}
[me@linuxbox ~]$ find playground \( -type f -not -perm 0600 \) -or \( -type d -not -perm 0700 \)
\end{lstlisting}
\end{latin}

این دستور تمامی ۲۶۰۰ فایل و ۱۰۰ دایرکتوری موجود (به اضافه‌ی فایل مرجع و خودِ دایرکتوری اصلی) را فهرست می‌کند، چرا که هیچ‌کدام با استانداردهای امنیتی فرضی ما (\cmd{0600} برای فایل‌ها و \cmd{0700} برای دایرکتوری‌ها) مطابقت ندارند. با تلفیق دانش خود پیرامون عملگرهای منطقی و کنش‌ها، می‌توانیم این دستور را به یک ابزار اصلاحی تبدیل کنیم تا سطوح دسترسی صحیح را به‌طور خودکار اعمال کند:

\begin{latin}
\begin{lstlisting}
[me@linuxbox ~]$ find playground \( -type f -not -perm 0600 -exec chmod 0600 '{}' ';' \) -or \( -type d -not -perm 0700 -exec chmod 0700 '{}' ';' \)
\end{lstlisting}
\end{latin}

در فعالیت‌های روزمره، شاید اجرای دو دستور مجزا برای فایل‌ها و دایرکتوری‌ها ساده‌تر به نظر برسد، اما درک این موضوع که چگونه می‌توان عملگرها و کنش‌ها را برای اجرای عملیات‌های پیچیده و مشروط ترکیب کرد، بسیار حائز اهمیت است.

\section{گزینه‌های کنترلی (Options) در دستور find}

در نهایت، به بررسی گزینه‌ها (\en{Options}) می‌پردازیم. گزینه‌ها در دستور \cmd{find} برای تعیین دامنه و محدوده‌ی جست‌وجو و همچنین کنترل رفتار پیمایشی ابزار به کار می‌روند. این گزینه‌ها را می‌توان در کنار سایر شرط‌ها و کنش‌ها برای دقیق‌تر کردن عبارات منطقی استفاده کرد. جدول زیر پرکاربردترین گزینه‌های \cmd{find} را معرفی می‌کند:

\begin{table}[htbp]
\centering
\caption{گزینه‌های کاربردی دستور \cmd{find}}
\label{tab:find-options}
\begin{tabularx}{\textwidth}{l X}
\toprule
\textbf{گزینه} & \textbf{شرح عملکرد فنی} \\
\midrule
\cmd{-depth} & به \cmd{find} دستور می‌دهد که هنگام پردازش یک دایرکتوری، ابتدا به سراغ محتویات آن (فایل‌ها و زیرشاخه‌ها) رفته و آن‌ها را پردازش کند و تنها پس از اتمام کار با محتویات، خودِ آن دایرکتوری را پردازش نماید. این شیوه پردازش به‌ویژه زمانی اهمیت پیدا می‌کند که قصد حذف یک دایرکتوری را دارید؛ زیرا دایرکتوری تنها زمانی قابل حذف است که کاملاً خالی باشد. کنش \cmd{-delete} به‌طور خودکار از این شیوه پردازش استفاده می‌کند تا فرآیند حذف دایرکتوری‌ها به‌درستی انجام شود. \\
\addlinespace
\cmd{-maxdepth levels} & حداکثر تعداد سطوح زیرشاخه‌ای که \cmd{find} باید در آن‌ها به جستجو بپردازد را مشخص می‌کند. به‌عنوان مثال، \cmd{-maxdepth 1} جستجو را به خودِ دایرکتوری شروع‌کننده محدود می‌کند و از ورود به زیرشاخه‌های آن جلوگیری می‌نماید. \\
\addlinespace
\cmd{-mindepth levels} & حداقل تعداد سطوحی را تعیین می‌کند که \cmd{find} باید پیش از اعمال شرط‌ها و کنش‌ها به آن‌ها نفوذ کند. برای نمونه، \cmd{-mindepth 1} باعث می‌شود که فایل‌های موجود در خودِ دایرکتوری شروع‌کننده (سطح صفر) نادیده گرفته شوند و جستجو از زیرشاخه‌های سطح اول آغاز گردد. \\
\addlinespace
\cmd{-mount} & (که با نام \cmd{-xdev} نیز شناخته می‌شود) از ورود \cmd{find} به دایرکتوری‌هایی که بر روی سیستم‌های فایل دیگری متصل (\en{mount}) شده‌اند، جلوگیری می‌کند. این گزینه جستجو را به همان سیستم فایل مبدأ محدود می‌کند و از پیمایش پارتیشن‌های دیگر جلوگیری می‌نماید. \\
\addlinespace
\cmd{-noleaf} & به \cmd{find} دستور می‌دهد که از بهینه‌سازی‌های ویژه‌ی سیستم‌های فایل یونیکسی (که بر اساس شمارش پیوندهای سخت (\en{hard links}) در دایرکتوری‌ها عمل می‌کنند) صرف‌نظر کند. این گزینه هنگام پیمایش سیستم‌های فایلی که از این ساختار پیروی نمی‌کنند، مانند رسانه‌های \en{CD-ROM} یا سیستم‌های فایل ویندوز (\en{FAT/NTFS})، ضروری است. \\
\bottomrule
\end{tabularx}
\end{table}

\section{جمع‌بندی}

با بررسی ویژگی‌های این دو ابزار، به وضوح مشخص است که دستور \cmd{locate} در مقایسه با ساختار لایه‌مند و چندوجهی \cmd{find}، سادگی بیشتری دارد. با این حال، هر یک از این دو ابزار در سناریوهای متفاوتی کاربرد تخصصی خود را دارند: \cmd{locate} برای جست‌وجوهای سریع در کل سیستم و \cmd{find} برای پایش دقیق و انجام عملیات‌های مشروط بر روی فایل‌ها.

توصیه می‌شود زمانی را به اکتشاف و تجربه در قابلیت‌های متعدد \cmd{find} اختصاص دهید. استمرار در به‌کارگیری این ابزار، علاوه بر افزایش بهره‌وری، درک عمیق‌تری از سازوکارهای مدیریت سیستم فایل در لینوکس به شما خواهد بخشید.

\section{منابع مطالعه تکمیلی}

برنامه‌های \cmd{locate}، \cmd{updatedb}، \cmd{find} و \cmd{xargs} همگی بخشی از بسته \cmd{findutils} پروژه گنو (\en{GNU Project}) هستند. این پروژه مستندات آنلاین گسترده و بسیار باکیفیتی را ارائه می‌دهد که مطالعه آن برای کاربرانی که از این ابزارها در محیط‌هایی با حساسیت امنیتی بالا استفاده می‌کنند، توصیه می‌شود:

\begin{latin}
\url{http://www.gnu.org/software/findutils/}
\end{latin}